\documentclass{article}
\usepackage[T1]{fontenc}
\usepackage{lmodern}
\usepackage{graphicx}
\usepackage{amsmath,amssymb}
\usepackage[a4paper,margin=2cm]{geometry}
\usepackage[absolute,overlay]{textpos}
\setlength{\TPHorizModule}{1mm}
\setlength{\TPVertModule}{1mm}
\usepackage[indonesian]{babel}
\usepackage{hyperref}
\setlength{\emergencystretch}{2em}
% unit: Halaman 10

\begin{document}

% === Halaman 10 (llm) ===

\section*{Dekripsi}

\begin{enumerate}
    \item Hitung kembali blok plainteks $m$ dari cipherteks $c$ menggunakan kunci privat $d$ dengan persamaan
    \[ m = c^d \pmod{n}, \]
    Jika cipherteks adalah blok $c_1, c_2, c_3, \dots$, dekripsi setiap blok menjadi
    \[ m_i = c_i^d \pmod{n} \]
    
    \item \textbf{Postprocessing}: Kodekan kembali $m$ atau $m_i$ menjadi bentuk pesan semula
\end{enumerate}

\end{document}
